\subsection{Symmetric tensors and Banach spaces}

13.1.1. Let E be a module over a commutative ring A. If n is an integer \(\geqslant0\), we denote by \(\mathsf{TS}^{n}(\mathbf{E})\) the module of symmetric tensors of degree n of E (A, III, p. 71 and A, IV, § 5, no. 2); the direct sum of the \(\mathsf{TS}^{n}(\mathbf{E})\) is denoted \(\mathsf{TS}(\mathbf{E})\). If r is either an integer or one of the symbols \(\infty\), \(\omega\), we set

\[
\mathsf{TS}^{(r)}(\mathrm{E})=\bigoplus_{n\leqslant r}\mathsf{TS}^{n}(\mathrm{E})\quad\text{and}\quad\mathsf{TS}^{(r)+}(\mathrm{E})=\bigoplus_{1\leqslant n\leqslant r}\mathsf{TS}^{n}(\mathrm{E});
\]

consequently, if \(r \geqslant 0\), \(\mathsf{TS}^{(r)}(\mathbf{E})\) is the direct sum of \(\mathsf{TS}^{(0)}(\mathbf{E}) = \mathbf{A}\) and \(\mathsf{TS}^{(r)+}(\mathbf{E})\). One has \(\mathsf{TS}^{(1)+}(\mathbf{E}) = \mathbf{E}\) and \(\mathsf{TS}^{(\infty)}(\mathbf{E}) = \mathsf{TS}^{(\omega)}(\mathbf{E}) = \mathsf{TS}(\mathbf{E})\).

If n is an integer \(\geqslant0\) and if \(x\in E\), one denotes by \(\gamma_{n}(x)\) the element \(x\otimes\cdots\otimes x\) of \(\mathsf{TS}^{n}(\mathbf{E})\); for \(n=0\), \(\gamma_{0}(x)=1\).

13.1.2 ("Polynomial maps"). Assume that A is an infinite integral domain and that E is a free A-module. Let F be an A-module. A map \(f\colon\mathbf{E}\to\mathbf{F}\) is called a homogeneous polynomial of degree \(n\) if it satisfies the following equivalent conditions (cf. A, IV, § 5, no. 9):

\begin{enumerate}
\def\labelenumi{\alph{enumi})}
\item
  There exists a linear map \(\tilde{f}\colon\mathsf{TS}^{n}(\mathbf{E})\to\mathbf{F}\) such that \(f(x)=\tilde{f}(\gamma_{n}(x))\) for all \(x\in\mathbf{E}\).
\item
  There exists a multilinear map \(u: \mathbf{E}^{n} \to \mathbf{F}\) such that \(f(x) = u(x, \ldots, x)\) for all \(x \in \mathbf{E}\).
\end{enumerate}

Suppose that this is the case. The map \(\tilde{f}\) satisfying \(a)\) is then unique; if \(t \in \mathsf{TS}^n(\mathbf{E})\), one denotes by \(\langle f, t \rangle\) the element \(\tilde{f}(t)\). If \(u: \mathbf{E}^n \to \mathbf{F}\) satisfies \(b)\), the corresponding linear map from \(\otimes^n\mathbf{E}\) to \(\mathbf{F}\) coincides with \(\tilde{f}\) on the submodule \(\mathsf{TS}^n(\mathbf{E})\). When \(n!\) is invertible in \(\mathbf{A}\), one can choose u symmetric, and do so uniquely.

13.1.3. Let E be a Banach space over K and F a separated polynormed space over K. Let U be an open neighborhood of 0 in E, \(f\colon\mathbf{U}\to\mathbf{F}\) a map of class \(\mathbf{C}^{r}\) (\(r\in\mathbf{N}_{\mathbb{K}}\)), and \(t\) an element of \(\mathsf{TS}^{(r)}(\mathbf{E})\), cf. 13.1.1 (applied to the ring A=K). Let \(k\) be an integer \(\leqslant r\) such that \(t\in\mathsf{TS}^{(k)}(\mathbf{E})\). One can decompose \(t\) and \(f\) in a unique way into

\[
t=t_{0}+\dots+t_{k},\quad\text{with}\quad t_{i}\in\mathrm{TS}^{i}(\mathrm{E})
\]

and

\[
f = f _ {0} + \dots + f _ {k} + h,
\]

where \(f_i\) is a continuous homogeneous polynomial of degree \(i\) ( \(0 \leqslant i \leqslant k\) ) and where \(h\) has contact of order \(\geqslant k\) with 0 at the point 0. One then sets

\[
\langle f, t \rangle = \sum_ {i = 0} ^ {k} \langle f _ {i}, t _ {i} \rangle ;
\]

this is an element of F that does not depend on the choice of k.

13.1.4. Let E and E' be two Banach spaces, U an open neighborhood of 0 in E, and \(\varphi\colon\mathbf{U}\to\mathbf{E}'\) a map of class \(\mathbf{C}^{r}\) such that \(\varphi(0)=0\). Let \(t\in\mathsf{TS}^{(r)}(\mathbf{E})\). There exists one and only one element \(t'\in\mathsf{TS}^{(r)}(\mathbf{E}')\) such that, for any separated polynormed space F, any open neighborhood \(\mathbf{U}'\) of 0 in \(\mathbf{E}'\), and any map \(f\colon\mathbf{U}'\to\mathbf{F}\) of class \(\mathbf{C}^{r}\), one has

\[
\langle f, t ^ {\prime} \rangle = \langle f \circ \varphi , t \rangle .\tag{*}
\]

This element \(t'\) is denoted \(\varphi_{*}(t)\). The map

\[
\varphi_ {*} \colon \mathsf {T S} ^ {(r)} (\mathrm{E}) \to \mathsf {T S} ^ {(r)} (\mathrm{E} ^ {\prime})
\]

thus defined is linear.
% semantic-fragments: legacy-input-seam
\relax{}
